\begin{exercice}\label{exo3-12} (\mortelexo)

Dans un cachot du Moyen-Age, le ge\^olier informe trois prisonniers que l'un d'entre eux sera choisi au hasard pour \^etre ex\'ecut\'e, alors que les deux autres seront lib\'er\'es. Le premier prisonnier demande alors au ge\^olier de lui dire quel prisonnier parmi les deux autres sera lib\'er\'e. Il se justifie en disant que de toutes fa\c{c}ons, il sait que l'un des deux sera libre, et que cette information ne change rien pour lui. 

\smallskip

Le ge\^olier lui r\'epond alors que s'il lui donne cette information, sa probabilit\'e d'\^etre ex\'ecut\'e pourrait augmenter, et qu'il pr\'ef\`ere alors garder le silence. %n'a donc pas int\'er\^et \`a conna\^itre ladite information... 

\smallskip

Que doit penser le prisonnier de la r\'eponse du ge\^olier, et de son propre sort ?

\medskip

%\corrref{TD1_1}

\end{exercice}
